
\documentclass{article}
\usepackage{colortbl}
\usepackage{makecell}
\usepackage{multirow}
\usepackage{supertabular}

\begin{document}

\newcounter{utterance}

\centering \large Interaction Transcript for game `cladder', experiment `full\_v1.5\_default', episode 10063 with Qwen3.6{-}27B.
\vspace{24pt}

{ \footnotesize  \setcounter{utterance}{1}
\setlength{\tabcolsep}{0pt}
\begin{supertabular}{c@{$\;$}|p{.15\linewidth}@{}p{.15\linewidth}p{.15\linewidth}p{.15\linewidth}p{.15\linewidth}p{.15\linewidth}}
    \# & \multicolumn{2}{c}{Player} && \multicolumn{2}{c}{Game Master} \\
    \hline

    \theutterance \stepcounter{utterance}  
    & & & \multicolumn{4}{p{0.6\linewidth}}{
        \cellcolor[rgb]{0.9,0.9,0.9}{
            \makecell[{{p{\linewidth}}}]{
                \texttt{\tiny{[P1$\langle$GM]}}
                \texttt{You are an expert in causal inference. The following question is not a typical commonsense query, but rather a meticulously designed question created by a professor specializing in causal inference, intended to assess the students' mastery of the course content.} \\
\texttt{Imagine a self{-}contained, hypothetical world with only the following conditions, and without any unmentioned factors or causal relationships: Zory has a direct effect on zuph and glimx. Zuph has a direct effect on jyka. Jyka has a direct effect on glimx. Zory is unobserved. We know that zory causes zuph. zuph causes not jyka. zory or jyka causes glimx. We observed an individual is zory. Would an individual is not glimx if not zuph instead of zuph?} \\
\texttt{Start your answer with 'Yes' or 'No', followed by additional reasoning or evidence to support your explanation.} \\
            }
        }
    }
     \\ \\

    \theutterance \stepcounter{utterance}  
    & \multicolumn{4}{p{0.6\linewidth}}{
        \cellcolor[rgb]{0.561,0.737,0.561}{
            \makecell[{{p{\linewidth}}}]{
                \texttt{\tiny{[P1$\rangle$GM]}}
                \texttt{No} \\
\\ 
\texttt{**Reasoning:**} \\
\\ 
\texttt{To determine the answer, we must reconstruct the causal structure and logical rules defined in the hypothetical world, then apply counterfactual reasoning.} \\
\\ 
\texttt{**1. Define the Causal Structure and Logical Rules:**} \\
\texttt{Based on the prompt, we have the following variables and relationships:} \\
\texttt{*   **Variables:** \$Zory\$ (\$Z\$), \$Zuph\$ (\$U\$), \$Jyka\$ (\$J\$), \$Glimx\$ (\$G\$).} \\
\texttt{*   **Causal Links:**} \\
\texttt{    1.  \$Z \textbackslash{}rightarrow U\$ (Zory causes Zuph)} \\
\texttt{    2.  \$U \textbackslash{}rightarrow \textbackslash{}neg J\$ (Zuph causes *not* Jyka. This implies if \$U\$ is true, \$J\$ is false. If \$U\$ is false, the effect on \$J\$ is not explicitly defined by this link alone, but we must look at other constraints. However, usually in these binary logic puzzles, "causes not J" implies \$U \textbackslash{}implies \textbackslash{}neg J\$. We need to determine the state of \$J\$ when \$U\$ is absent. Let's look at the next rule.)} \\
\texttt{    3.  \$Z \textbackslash{}rightarrow G\$ (Zory has a direct effect on Glimx)} \\
\texttt{    4.  \$J \textbackslash{}rightarrow G\$ (Jyka has a direct effect on Glimx)} \\
            }
        }
    }
    & & \\ \\

    \theutterance \stepcounter{utterance}  
    & & & \multicolumn{2}{p{0.3\linewidth}}{
        \cellcolor[rgb]{0.9,0.9,0.9}{
            \makecell[{{p{\linewidth}}}]{
                \texttt{\tiny{[GM$|$GM]}}
                \texttt{No} \\
            }
        }
    }
    & & \\ \\

    \theutterance \stepcounter{utterance}  
    & & & \multicolumn{2}{p{0.3\linewidth}}{
        \cellcolor[rgb]{0.9,0.9,0.9}{
            \makecell[{{p{\linewidth}}}]{
                \texttt{\tiny{[GM$|$GM]}}
                \texttt{no} \\
            }
        }
    }
    & & \\ \\

    \theutterance \stepcounter{utterance}  
    & & & \multicolumn{2}{p{0.3\linewidth}}{
        \cellcolor[rgb]{0.9,0.9,0.9}{
            \makecell[{{p{\linewidth}}}]{
                \texttt{\tiny{[GM$|$GM]}}
                \texttt{game\_result = WIN} \\
            }
        }
    }
    & & \\ \\

\end{supertabular}
}

\end{document}
